\subsection{Semi-simple and absolutely semi-simple endomorphisms}

DEFINITION 6. --- Let E be a finite dimensional vector space over a commutative field K. Then an endomorphism u of E is said to be semi-simple if every subspace of E which is closed under u has a complement which is closed under u.

This means that every submodule of the K{[}X{]}-module \(E_{u}\) is a direct factor, in other words the K{[}X{]}-module \(E_{u}\) is semi-simple (VII, p.~9).

PROPOSITION 14. --- An endomorphism u of a finite dimensional vector space over a commutative field is semi-simple if and only if the minimal polynomial of u has no multiple factors.

This follows immediately from VII, p.~9, Cor. 4 and p.~31, Remark 1.

Let E be a vector space over a commutative field K, let L be an extension of K, and u an endomorphism of E; let \(u_{(L)}\) denote the L-endomorphism \(1_{L} \otimes u\) of the L-vector space \(\mathrm{E}_{(L)} = \mathrm{L} \otimes_{K} \mathrm{E}\) induced from E by extension of scalars. In the same way, if F is a set of endomorphisms of E, let \(\mathfrak{F}_{(L)}\) denote the set of \(u_{(L)}\) for u in F.

COROLLARY. --- Let u be an endomorphism of a finite dimensional vector space over a commutative field K, and let L be an extension of K. If \(u_{(L)}\) is semi-simple then u is semi-simple. If u is semi-simple and L is separable over K, then \(u_{(L)}\) is semi-simple.

This follows immediately from Prop. 14 and from V, p.~120, Cor. 1 (note that the minimal polynomials of u and \(u_{(L)}\) coincide).

PROPOSITION 15. --- Let E be a finite dimensional vector space over a commutative field K, let u be an endomorphism of E and let \(q(X)\) be its minimal polynomial. Then the following conditions are equivalent:

\begin{enumerate}
\def\labelenumi{(\roman{enumi})}
\item
  For every extension \(\mathbf{L}\) of \(\mathbf{K}\) , the endomorphism \(u_{(\mathbf{L})}\) is semi-simple.
\item
  There exists an extension \(\mathbf{L}\) of \(\mathbf{K}\) such that the endomorphism \(u_{(\mathbf{L})}\) is diagonalisable.
\item
  The polynomial \(q(\mathbf{X})\) is separable over \(\mathbf{K}\) .
\end{enumerate}

Indeed, condition (i) means that the polynomial \(1 \otimes q(\mathbf{X})\) in L{[}X{]} has no multiple factors, for any extension L of K (Prop. 14), condition (ii) means that there exists an extension L of K such that all the roots of \(q(\mathbf{X})\) belong to L and that these roots are all simple (VII, p.~40, Prop. 12), and these conditions are each equivalent to (iii) by definition (V, p.~38).

DEFINITION 7. --- An endomorphism u satisfying conditions (i), (ii) and (iii) of Prop. 15 is said to be absolutely semi-simple.

COROLLARY. --- A necessary and sufficient condition for u to be absolutely semi-simple is that there exist an extension L of K such that L is perfect and \(u_{(L)}\) is semi-simple.

The condition in the corollary means that there exists an extension L of K such that L is perfect and \(q(\mathbf{X})\) has no multiple factors in L{[}X{]} (Prop. 14); this condition is equivalent to (iii) by V, p.~38, Cor. 2.

PROPOSITION 16. --- Let E be a finite dimensional vector space over a commutative field K, let F be a set of endomorphisms of E and let A be the subalgebra of \(\mathrm{End}_{\mathbf{K}}(\mathbf{E})\) generated by F and \(\mathrm{Id}_{\mathbf{E}}\) . Then the following conditions are equivalent:

\begin{enumerate}
\def\labelenumi{(\roman{enumi})}
\item
  There exists an extension \(\mathbf{L}\) of \(\mathbf{K}\) such that \(\mathfrak{F}_{(\mathbb{L})}\) is diagonalisable.
\item
  The K-algebra A is étale (V, p.~28, Def. 1).
\item
  The elements of \(\mathfrak{F}\) are absolutely semi-simple and commute with one another.
\end{enumerate}

Note first that, for every extension L of K, the L-algebra generated by \(\mathfrak{F}_{(L)}\) and \(Id_{E_{(L)}}\) coincides with \(L \otimes_{K} A\) ; hence by Prop. 13 \(\mathfrak{F}_{(L)}\) is diagonalisable if and only if the L-algebra \(L \otimes_{K} A\) is diagonalisable. The equivalence of conditions (i) and (ii) thus follows from V, p.~28, Def. 1. On the other hand, it is immediate that (i) \(\Rightarrow\) (iii). Finally suppose (iii) holds, and let L be an algebraic closure of K; then the elements of \(\mathfrak{F}_{(L)}\) are diagonalisable (VII, p.~40, Prop. 12) and commute with one another; hence \(\mathfrak{F}_{(L)}\) is diagonalisable by VII, p.~41, Prop. 13.

COROLLARY. --- The sum and product of two commuting, absolutely semi-simple endomorphisms are absolutely semi-simple.

Remark. --- Suppose the conditions of Prop. 16 are satisfied and let L be an extension of K. By Prop. 13 the set \(\mathfrak{F}_{(L)}\) is diagonalisable if and only if the algebra \(L \otimes_K A\) is diagonalisable. It follows from V, p.~30, Prop. 2 that there exists a finite extension L of K such that \(\mathfrak{F}_{(L)}\) is diagonalisable. In fact L may be taken to be Galois; indeed, taking a finite subset \(\mathfrak{F}'\) of \(\mathfrak{F}\) which generates A, we may take L to be a splitting field for the minimal polynomials of the elements of \(\mathfrak{F}'\) (Prop. 12 and 13).

